% s-178-p178-m001; original PDF 178
\Aphorism{s-178-p178-m001}{266}
Ich blicke von einem schroffen Felsencap hinunter in die enge Bucht, in der das Hafenstädtchen liegt. Es ist ein sonnenloser Spätnachmittag und jene Stille, die nichts Todtes, sondern nur etwas grenzenlos Versunkenes und Nachdenkliches hat: das halbtraumhafte Sinnen eines Einsiedlers, dem derweilen der Bart über die Füsse wächst. Und doch auch etwas Rathloses und Leeres — man sieht kein Ende des Grübelns und Träumens. Weit hinaus ins Meer, das leise und sehr langsam Athem holt, hängen trübe Wolken, die doch nicht regnen wollen; nur ganz in der Ferne ist freier Himmel und ein sonnenbeschienener Streifen Wassers. Aber was sehe ich! da draussen am Horizont schwebt eine Lichtsäule; es muss wohl ein Segel sein, auf das späte Sonne glüht. Dieser leuchtende Schmetterling, der schweigsam und braungoldenen Scheines vorüberzieht, ist das Schwermüthigste in der schwermüthigen Stille. Ich denke dabei an eine einsam schweifende Seele, die immer fern von bewohnten Küsten hingleiten muss; dafür fängt sie allein den Schimmer der sinkenden Sonne auf und glänzt landeinwärts, tief hinein in die Bucht, unter hängende Wolken und in beschattete Häfen. Mich dünkt, ich höre die Klage dieser irrenden Seele in langgezogenen Tönen über das Wasser streichen. „Ach, dass ich dem verworrenen Schattenspiel Leuchte sein muss! dass ich hinaus soll in diese brandende Meerbreite, um trüben dämmernden Menschenherzen einen Sonnenabglanz vorzugaukeln! Nun dunkelt der Abend, und im Hafen wird Ruh; wer noch Licht will, entzündet sein Öllämpchen — nur ich muss unerbittlich dem Urlicht nach, ich selig-unseligstes der Segel!“
\par

